\subsection{A sharp normalization obstruction}
\label{sec:rms-obstruction}

Theorem~\ref{thm:main-rms} assumes a certified denominator floor.
This subsection shows that the floor matters: after uniform attention,
an unstabilized normalization can force a linear number of context
probes, and a stabilizer $\epsilon$ caps the cost at order
$\epsilon^{-1}$. Corollary~\ref{cor:rmsfixedgrid} keeps every weight
on one fixed grid.

\begin{theorem}[A sharp normalization family]
\label{thm:rmscontextlaw}
For every $T\ge2$ and rational $0\le\epsilon\le1$, there is a
width-two causal block consisting of uniform attention, RMSNorm,
and a tanh map, each replacing its input stream, with row norms at
most one and two logits
in $[-1,1]$, on the two-token embedding alphabet
$\{(-1,1),(1,1)\}$, such that its exact-sampling query optimum satisfies
\begin{equation}
 \frac1{384}\min\{T,\epsilon^{-1}\}
 \ \le\ Q^*\ \le\
 64\min\{T,\epsilon^{-1}\}.
 \label{eq:rmscontextlaw}
\end{equation}
Here $\epsilon^{-1}=+\infty$ at zero. The lower bound applies to all
exact samplers, with arbitrary weight-only preprocessing. The upper
bound has an implementation using
$O(B_N^4\min\{T,\epsilon^{-1}\})$ expected bit operations.
All nonzero network coefficients are dyadic with
$O(\log T)$ bits. The denominator is nonzero on every context,
including when $\epsilon=0$.
\end{theorem}

The score matrices in this construction are zero. Consequently the
lower bound holds at every inverse temperature, including $\beta=0$.

The matching sampler uses the following identity.

\begin{lemma}[A geometric majority identity]
\label{lem:rmsgeommajority}
Let $\theta>0$. Draw $K\ge0$ with
\[
 \Pr(K=k)=\frac\theta{1+\theta}(1+\theta)^{-k}.
\]
The majority of $2K+1$ independent signs of mean $z$ has mean
\begin{equation}
 G_\theta(z)=\frac{z\sqrt{1+\theta}}{\sqrt{\theta+z^2}}.
 \label{eq:rmsgeommajority}
\end{equation}
Its expected number of source signs is $1+2/\theta$.
\end{lemma}
\paragraph{Proof idea.}
Differentiate the majority polynomials before summing the geometric
mixture. Their derivatives form a binomial series with exponent $-3/2$,
which is also the derivative of the required normalized mean.

\begin{proof}
Write $M_k(z)$ for the mean of majority of $2k+1$ signs. The proof
of Lemma~\ref{lem:comptanharity} gives $M_k'(z)=C_k(1-z^2)^k$ with
$C_k=(2k+1)\binom{2k}{k}4^{-k}$, and $C_k=(3/2)_k/k!$.
The geometric mixture can be differentiated term by term
uniformly on $[-1,1]$. The binomial series gives its derivative
as
\[
 \frac\theta{1+\theta}
 \left(1-\frac{1-z^2}{1+\theta}\right)^{-3/2}
 =\frac{\theta\sqrt{1+\theta}}{(\theta+z^2)^{3/2}}
 =G_\theta'(z).
\]
Both functions vanish at zero, proving the identity. The
arity expectation follows from $\mathbb E K=1/\theta$.
\end{proof}

\paragraph{Proof idea for Theorem~\ref{thm:rmscontextlaw}.}
Uniform attention writes the context mean $M$ next to a small dyadic
sentinel, and normalization makes the output a steep function of $M$
near zero. The transcript score turns this steepness into order
$\min\{T,\epsilon^{-1}\}$ probes. A geometric mixture of majorities
realizes the same function at that cost, and a full scan is cheaper
when $\epsilon<1/T$.

\begin{proof}[Proof of Theorem~\ref{thm:rmscontextlaw}]
Take width two and contexts
\[
 h_{0,j}=(x_j,1),\qquad x_j\in\{-1,+1\}.
\]
Let $c=2^{-\lceil\log_2T\rceil}$, so $0<c\le1/T$. The
query and key maps are zero, and the value map is
$\operatorname{diag}(1,c)$. At the final position, uniform
causal attention gives
\[
 a=(M,c),\qquad M=T^{-1}\sum_{j=1}^T x_j.
\]
Apply \eqref{eq:rmsdefinition}, then a tanh row with coefficient
$1/2$ on the first coordinate. Give the two output categories
opposite logits. If the category is encoded by $Y\in\{-1,+1\}$,
its target mean is
\begin{equation}
 f_\epsilon(M)
 =\tanh\!\left(\tanh\!\left(
   \frac{M}{\sqrt2\sqrt{2\epsilon+c^2+M^2}}
 \right)\right).
 \label{eq:rmsfamilymean}
\end{equation}
Every parameter except $\epsilon$ is dyadic. The value and
tanh row norms are at most one, and the logits lie in $[-1,1]$.
The sentinel $c$ prevents a zero denominator.

\paragraph{The lower bound.}
Lemma~\ref{lem:transcript} gives $H'(0)^2\le\mathbb E_0Q$ by
Cauchy--Schwarz. This is the finite-input form of the sequential
information bound; the first-score Bernoulli-factory lower bound
is standard \cite{mendo2019}.

For $0\le u\le1$, $\tanh u\ge u/2$. The argument of the
inner tanh in \eqref{eq:rmsfamilymean} has absolute value
at most $1/\sqrt2$. Oddness consequently gives
\[
 M f_\epsilon(M)
 \ge\frac{M^2}{8\sqrt{2\epsilon+c^2+M^2}}.
\]
For independent fair signs,
\[
 \mathbb E M^2=\frac1T,
 \qquad \mathbb E M^4=\frac{3T-2}{T^3}\le\frac3{T^2}.
\]
For any nonnegative $Z$ with positive mean and any $a>0$,
weighted Jensen gives
\begin{equation}
 \mathbb E\frac{Z}{\sqrt{a+Z}}
 \ge\frac{\mathbb E Z}
 {\sqrt{a+\mathbb E Z^2/\mathbb E Z}}.
 \label{eq:rmsweightedjensen}
\end{equation}
To see this, tilt the probability law by $Z/\mathbb E Z$
and apply convexity of $u\mapsto(a+u)^{-1/2}$.
Taking $Z=M^2$ yields
\[
 H'(0)=T\mathbb E[Mf_\epsilon(M)]
 \ge\frac1{8\sqrt{2\epsilon+c^2+3/T}}.
\]
Lemma~\ref{lem:transcript} now implies
\begin{equation}
 \mathbb E_0Q
 \ge\frac1{64(2\epsilon+c^2+3/T)}
 \ge\frac1{384}\min\{T,\epsilon^{-1}\}.
 \label{eq:rmsscorelower}
\end{equation}
The last inequality uses $c^2\le1/T^2\le1/T$ and considers
$\epsilon\le1/T$ and $\epsilon\ge1/T$ separately.
An average-case lower bound is also a worst-context lower
bound. This proves the lower part of
Theorem~\ref{thm:rmscontextlaw}.

\paragraph{A matching sampler.}
For \eqref{eq:rmsfamilymean}, take $\theta=2\epsilon+c^2$.
A known gate of probability
$[2(1+\theta)]^{-1/2}$ turns the majority sign into a sign of
mean $M/[\sqrt2\sqrt{\theta+M^2}]$, using a fair sign on the
other branch. A source sign of mean $M$ is obtained by querying
one uniform context position.

Lemma~\ref{lem:unittanhfour} supplies a sign of mean $\tanh z$
from signs of mean $z\in[-1,1]$. Its expected source count is at
most two, which is in particular at most $2e$, and its local bit
work is $O(B_N^4)$.

Applying this tanh factory twice realizes
\eqref{eq:rmsfamilymean}. Its expected number of context
requests is at most
\[
 4e^2(1+2/\theta)
 \le32(1+\epsilon^{-1})
 \le64\epsilon^{-1},\qquad 0<\epsilon\le1.
\]
Alternatively, read all $T$ inputs, form the exact rational
$M$, and use certified square roots and exponentials to
sample the two-category probability. Choose this alternative
when $\epsilon<1/T$. The first construction is used when
$\epsilon\ge1/T$. Every geometric parameter then has inverse
at most $1+T$, and its arity and arithmetic tails give
$O(B_N^4/\epsilon)$ expected bit work. In the scan case,
the sentinel is at least $1/(2T)$, so the additional precision
needed by the square root is $O(\log T)$ and the total work
is $O(TB_N^4)$. This proves both upper bounds in
Theorem~\ref{thm:rmscontextlaw}.
\end{proof}

\subsubsection{A fixed coefficient grid}

The width-two construction uses a weight with $O(\log T)$ bits.
That precision is permitted by the finite-encoding model above,
but a fixed grid such as $b=20\lceil\log_2n\rceil$
would not allow this weight at fixed width and unbounded $T$.
The next construction generates its sentinel through averaging.
Its weights are independent of the context length.

\begin{corollary}[A fixed-grid normalization law]
\label{cor:rmsfixedgrid}
There is a width-four block consisting of uniform attention, RMSNorm,
and tanh, each replacing its input stream, with row norms at most
one and nonzero matrix
coefficients belong to $\{1/2,1\}$, such that
\begin{equation}
 \frac1{128}\min\{T,\epsilon^{-1}\}
 \le Q^*\le32\min\{T,\epsilon^{-1}\},
 \qquad 0\le\epsilon\le1.
 \label{eq:rmsfixedgridlaw}
\end{equation}
The upper bound has
$O(B_N^4\min\{T,\epsilon^{-1}\})$ expected bit work. The
input has one unknown sign per position and a public
first-position feature. Every normalization denominator is
positive at every causal position on every allowed context.
In particular the unstabilized linear lower bound holds with
one fixed coefficient grid as $T$ tends to infinity.
\end{corollary}
\paragraph{Proof idea.}
A public first-position feature produces the sentinel $1/t$ through
uniform averaging, so no context-dependent matrix coefficient is needed.
The score lower bound and geometric-majority sampler then apply to the
resulting scalar normalization, with a full scan available at small
stabilization.

\begin{proof}
Use
\[
 h_{0,j}=(x_j,1,\eta_j,1),\qquad
 \eta_1=1,\quad \eta_j=-1\ (j>1).
\]
The first-position feature $\eta_j$ is public and requires no
probe of $x_j$. Both query and key maps vanish. The first
value row selects $x_j$, the second takes the average of
coordinates two and three, and the other rows vanish.
Consequently, at every causal position $t$,
\[
 y_t=(M_t,1/t,0,0),\qquad
 M_t=t^{-1}\sum_{j\le t}x_j.
\]
The sentinel is strictly positive for every $t$. The tanh
output row is one half of the first normalized coordinate.
Opposite output logits give target sign mean
\[
 \widetilde f_\epsilon(M)
 =\tanh\!\left(\tanh\!\left(
 \frac{M}{\sqrt{4\epsilon+T^{-2}+M^2}}
 \right)\right).
\]
All matrix rows have norm at most one and all biases are zero.
The inner argument belongs to $[-1,1]$. The same two tanh
bounds give
\[
 M\widetilde f_\epsilon(M)
 \ge\frac{M^2}{4\sqrt{4\epsilon+T^{-2}+M^2}}.
\]
Weighted Jensen and the finite-input score lemma imply
\[
 \mathbb E_0Q
 \ge\frac1{16(4\epsilon+T^{-2}+3/T)}
 \ge\frac1{128}\min\{T,\epsilon^{-1}\}.
\]
For the last step, use $T^{-2}\le1/T$ and
$\epsilon+1/T\le2\max\{\epsilon,1/T\}$.

For the upper bound, put $\theta=4\epsilon+T^{-2}$ in
Lemma~\ref{lem:rmsgeommajority}. A known gate of probability
$(1+\theta)^{-1/2}$ gives the inner argument. Applying the
unit-interval tanh factory from Lemma~\ref{lem:unittanhfour}
twice costs fewer than sixteen requests of that source sign.
Thus for $0<\epsilon\le1$ the expected number of context
requests is at most
\[
 16(1+2/\theta)
 \le16(1+(2\epsilon)^{-1})
 \le24\epsilon^{-1}.
\]
When $\epsilon<1/T$, scan the context instead. This gives
the constant 32 in \eqref{eq:rmsfixedgridlaw} with slack.
The bit implementation and its precision bound are identical
to those in Theorem~\ref{thm:rmscontextlaw}; now the rational
sentinel $1/T$ is computed from the position index, without
being a network coefficient.
\end{proof}
